\section{Trouver le MST d'un graphe}

Un arbre couvrant minimal (MST) est un arbre couvrant dont le poids total est le plus petit possible.

Dans le cadre du projet, trois algorithmes classiques de construction d'un MST ont été réalisés : l'algorithme de Kruskal, l'algorithme de Prim et l'algorithme par suppression inverse.

\subsection{Algorithme de Kruskal}
Kruskal trie les arêtes par poids et les ajoute une à une à l'arbre, en évitant de former des cycles.
La première version implémentait une vérification explicite des cycles, ce qui pour un graphe complet de 1000 sommets prenait environ une seconde.
La version améliorée utilise une structure DSU (Union-Find), ce qui réduit le temps de calcul d'environ moitié.
Cet algorithme s'est révélé très rapide par rapport aux autres méthodes.

\subsection{Algorithme de Prim}
Prim construit l'arbre à partir d'un sommet de départ : à chaque étape, on choisit l'arête de poids minimal qui relie l'arbre partiel à un sommet non encore inclus.
Dans notre implémentation, cette version a montré des performances bien inférieures à Kruskal.
Mais on peut l'améliorer considérablement en utilisant une structure de données différente pour stocker les arêtes.

\subsection{Algorithme par suppression inverse}
Cette méthode trie les arêtes par poids décroissant et tente de supprimer successivement l'arête la plus lourde tout en conservant la connexité.
Elle s'est avérée plus lente comparée àux l'autres.
On peut peut l'améliorer en améliorant la façon dont on vérifie que le graphe est connexe.
